%% kapa: exam=1 round=37 year=2026 period=1 form=A subject=civil_law subject_ko=민법 seq=31 number=31
%% source: Q-Net 감정평가사 기출문제 (cst00309) artlSeq=5258924 file=raw/1008/5258924/2026년 제37회 감정평가사 1차 1교시.pdf
%% conversion: kordoc 4.24.51 -> tools/gen1.py (auto); review=proofread-sonnet; flags=none
\begin{question}{31}
甲은 그 소유의 X토지 위에 Y건물을 신축하여 소유하고자 하는 乙과 X토지에 관하여 지상권 설정계약을 체결하였다. 이에 관한 설명으로 옳지 않은 것은? (다툼이 있으면 판례에 따름)
\begin{choices}
  \item X토지의 일부에 관하여도 범위를 특정하여 등기하면 지상권은 성립할 수 있다.
  \item 乙은 지상권을 유보한 채 신축된 Y건물의 소유권만을 양도할 수 있다.
  \item 甲과 乙의 지상권설정계약에서 지료의 지급에 대한 약정이 없었더라도 지상권의 성립에는 영향이 없다.
  \item 지상권의 존속기간이 만료되지 않았더라도 신축된 Y건물이 멸실되면 지상권은 소멸한다.
  \item 지상권이 소멸한 경우, 甲이 상당한 가액을 제공하여 현존하는 Y건물의 매수를 청구한 때에는 乙은 정당한 이유 없이 이를 거절하지 못한다.
\end{choices}
\end{question}
\answer{4}
